\subsection{Extension class associated with a regular sequence}

Let A be a ring, M an A-module, \(x = (x_{1}, \ldots, x_{n})\) a sequence of elements of A. Denote by \(M_{i}\) the A-module \(\mathbf{M}/(x_{1}\mathbf{M} + \cdots + x_{i-1}\mathbf{M})\) for \(i = 0, \ldots, n + 1\) , so that \(M_{0}\) and \(M_{1}\) identify with M and that \(\mathbf{M}_{n+1} = \mathbf{M}/(\mathbf{x})\mathbf{M}\) . Let

\[
\overline {{x}} _ {i}: \mathbf {M} _ {i - 1} \to \mathbf {M} _ {i}, \qquad i = 1, \dots , n,
\]

denote the A-homomorphism composed of the homothety of \(M_{i-1}\) of ratio \(x_{i}\) and of the canonical projection of \(M_{i-1}\) over \(M_{i}\) . Finally, let us denote by \(p: M_{n} \to M/(x)M\) the canonical projection. The diagram

\[
0 \longrightarrow \mathrm{M} \xrightarrow {\overline {{{x}}} _ {1}} \mathrm{M} _ {1} \xrightarrow {\overline {{{x}}} _ {2}} \mathrm{M} _ {2} \longrightarrow \dots \xrightarrow {\overline {{{x}}} _ {n}} \mathrm{M} _ {n} \xrightarrow {p} \mathrm{M} / (\boldsymbol {x}) \mathrm{M} \longrightarrow 0\tag{19}
\]

is an exact sequence if and only if the sequence x is M-regular. Suppose henceforth that the sequence x is regular for M. The element \(\theta_{x} \in \operatorname{Ext}_{A}^{n}(\mathbf{M}/(\mathbf{x})\mathbf{M}, \mathbf{M})\) associated with the exact sequence (19) is also said to be associated with the M-regular sequence x.

Let i be an integer, \(1 \leqslant i \leqslant n\) . Note that the sequence (19) can be decomposed into the two exact sequences

\begin{enumerate}
\def\labelenumi{(\arabic{enumi})}
\setcounter{enumi}{19}
\tightlist
\item
\end{enumerate}

\[
0 \longrightarrow \mathrm{M} \xrightarrow {\overline {{x}} _ {1}} \mathrm{M} _ {1} \xrightarrow {\overline {{x}} _ {2}} \mathrm{M} _ {2} \longrightarrow \dots \xrightarrow {\overline {{x}} _ {i}} \mathrm{M} _ {i} \longrightarrow \mathrm{M} / (x _ {1} \mathrm{M} + \dots + x _ {i} \mathrm{M}) \longrightarrow 0\tag{21}
\]

\[
0 \longrightarrow \mathrm{M} / (x _ {1} \mathrm{M} + \dots + x _ {i} \mathrm{M}) \xrightarrow {\overline {{x}} _ {i + 1}} \mathrm{M} _ {i + 1} \longrightarrow \dots
\]

\[
\longrightarrow \mathrm{M} _ {n} \stackrel {p} {\longrightarrow} \mathrm{M} / (x) \mathrm{M} \longrightarrow 0
\]

which are none other than the exact sequences associated with the sequence \((x_{1}, \ldots, x_{i})\) which is regular for M and to the sequence \((x_{i+1}, \ldots, x_{n})\) which is regular for \(\mathrm{M}/(x_{1} \mathrm{~M} + \cdots + x_{i} \mathrm{~M})\) . Denoting

\[
\theta_ {(x _ {1}, \dots , x _ {i})} \in \operatorname{Ext} _ {\mathbf {A}} ^ {i} (\mathbf {M} / (x _ {1} \mathbf {M} + \dots + x _ {i} \mathbf {M}), \mathbf {M})
\]

\[
\theta_ {(x _ {i + 1}, \dots , x _ {n})} \in \operatorname{Ext} _ {\mathbf {A}} ^ {n - i} (\mathbf {M} / (\boldsymbol {x}) \mathbf {M}, \mathbf {M} / (x _ {1} \mathbf {M} + \dots + x _ {i} \mathbf {M}))
\]

the extension classes associated with (20) and (21), we have, according to X, p.~118, prop. 3

\[
\theta_ {(x _ {1}, \dots , x _ {n})} = \theta_ {(x _ {1}, \dots , x _ {i})} \circ \theta_ {(x _ {i + 1}, \dots , x _ {n})}.\tag{22}
\]

Moreover, by prop. 5 (X, p.~157), the Koszul complex K.(x, M) is acyclic in degrees \(\neq n\) , whence an exact sequence

\[
0 \longrightarrow \mathrm{M} \xrightarrow {\partial^ {0}} \mathbf {K} ^ {1} (\boldsymbol {x}, \mathrm{M}) \xrightarrow {\partial^ {1}} \mathbf {K} ^ {2} (\boldsymbol {x}, \mathrm{M}) \xrightarrow {\partial^ {2}} \dots \longrightarrow \mathbf {K} ^ {n} (\boldsymbol {x}, \mathrm{M}) \xrightarrow {q} \mathrm{M} / (\boldsymbol {x}) \mathrm{M} \longrightarrow 0\tag{23}
\]

where we have identified \(\mathbf{K}^{0}(\boldsymbol{x},\mathbf{M})\) to \(\mathbf{M}\) and where \(q\) maps each element \(\boldsymbol{m}\in\mathbf{K}^{n}(\boldsymbol{x},\mathbf{M})\) on the class in \(\mathbf{M}/(\boldsymbol{x})\mathbf{M}\) of \(\boldsymbol{m}(1,2,\ldots,n)\in\mathbf{M}\) .

PROPOSITION 8. --- Suppose the sequence x is regular for M. The element of Ext \(_{A}^{n}\) (M/(x) M, M) associated with the exact sequence (23) is (-1) \(^{n(n+1)/2}\) . θ \(_{x}\) .

For \(i = 0,1,\dots,n\) let us define an A-linear map

\[
a ^ {i}: \mathbf {K} ^ {i} (\boldsymbol {x}, \mathrm{M}) \rightarrow \mathrm{M} _ {i} = \mathrm{M} / (x _ {1} \mathrm{M} + \dots + x _ {i - 1} \mathrm{M})
\]

as follows: if \(\boldsymbol{m} \in \mathbf{K}^i(\boldsymbol{x}, \mathbf{M})\) , \(a^i(\boldsymbol{m})\) is the class in \(\mathbf{M}_i\) of the element \(\boldsymbol{m}(1, 2, \ldots, i)\) of M. It is clear that \(a^0\) is the identity map of M and that \(p \circ a^n = q\) . Moreover \(a^{i+1} \circ \partial^i(\boldsymbol{m})\) is the image in \(\mathbf{M}_{i+1}\) of the element

\[
\sum_ {k = 1} ^ {i + 1} (- 1) ^ {k + 1} x _ {k} \boldsymbol {m} (1, 2, \dots , k - 1, k + 1, \dots , i + 1).
\]

Since \(x_{k}\) annihilates \(\mathbf{M}_{i+1}\) for \(k=1,\ldots,i,a^{i+1}\circ\partial^{i}(\boldsymbol{m})\) is also the image of

\[
(- 1) ^ {i} x _ {i + 1} \boldsymbol {m} (1, 2, \dots , i),
\]

hence

\[
a ^ {i + 1} \circ \partial^ {i} = (- 1) ^ {i} \overline {{{x}}} _ {i + 1} \circ a ^ {i}.
\]

By X, p.~120, cor. 1 and 2, the element of Ext \(_{A}^{n}\) (M/(x) M, M) associated with (23) is equal to \(\prod_{i=1}^{n} (-1)^{i} \cdot \theta_{x}\) , whence the assertion.

COROLLARY. — Suppose moreover that the modules \(\mathbf{M}/(x_{1}\mathbf{M}+\cdots+x_{i-1}\mathbf{M})\) are separated for the (x)-adic topology, and let \((a_{ij})\in\mathbf{GL}_{n}(\mathbf{A})\) . Set

\[
y _ {i} = \sum_ {j} a _ {i j} x _ {j} \text { and } \mathbf {y} = (y _ {1}, \dots , y _ {n}).
\]

Then the sequence y is regular for M, and one has \(\theta_{y} = \det (a_{ij})^{-1}\theta_{x}\) .

Indeed, the sequence y is regular for M by prop. 6 (X, p.~158) and theorem 1 (X, p.~160); the last assertion follows from prop. 8, and from prop. 4 of X, p.~119.

PROPOSITION 9. — Suppose the sequence x is regular for M. If N is an A-module such that (x) N = 0, one has Ext \(_{A}^{i}\) (N, M) = 0 for i \textless{} n, and the map α → θ \(_{x}\) ∘ α from Hom \(_{A}\) (N, M/(x) M) to Ext \(_{A}^{n}\) (N, M) (which is also the iterated connecting homomorphism associated with (19), cf.~X, p.~127, cor. 3) is bijective.

It remains to prove that the homomorphism \(\psi^{i}:\alpha\mapsto\theta_{x}\circ\alpha\) of \(\operatorname{Ext}_{\mathbf{A}}^{i-n}(\mathbf{N},\mathbf{M}/(\mathbf{x})\mathbf{M})\) to \(\operatorname{Ext}_{\mathbf{A}}^{i}(\mathbf{N},\mathbf{M})\) is bijective for \(i\leqslant n\) . Let us reason by induction on \(n\) , the assertion being trivial for \(n=0\) . Set \(\mathbf{M}_{1}=\mathbf{M}/x_{1}\mathbf{M}\) , \(\mathbf{x}^{\prime}=(x_{2},...,x_{n})\) , so that \(\mathbf{x}^{\prime}\) is \(\mathbf{M}_{1}\) -regular. By the induction hypothesis, the homomorphism

\[
\overline {{{{\psi}}}} ^ {i - 1}: \alpha \mapsto \theta_ {\boldsymbol {x} ^ {\prime}} \circ \alpha
\]

of \(\operatorname{Ext}_{\mathbf{A}}^{i-n}(\mathbf{N}, \mathbf{M}/(\mathbf{x})\mathbf{M})\) to \(\operatorname{Ext}_{\mathbf{A}}^{i-1}(\mathbf{N}, \mathbf{M}_1)\) is bijective for \(i < n\) . Moreover, consider the exact sequence

\[
0 \longrightarrow \mathrm{M} \xrightarrow {(x _ {1}) _ {\mathrm{M}}} \mathrm{M} \longrightarrow \mathrm{M} _ {1} \longrightarrow 0;
\]

the connecting homomorphism \(\operatorname{Ext}_{\mathbf{A}}^{i}(\mathbf{N}, \mathbf{M}_{1}) \to \operatorname{Ext}_{\mathbf{A}}^{i+1}(\mathbf{N}, \mathbf{M})\) associated is \(\varphi^{i}: \beta \mapsto \theta_{x_{1}} \circ \beta\) (X, p.~125, prop. 5);

since we have \(\operatorname{Ext}^i(1_{\mathbf{N}}, (x_1)_{\mathbf{M}}) = \operatorname{Ext}^i((x_1)_{\mathbf{N}}, 1_{\mathbf{M}}) = 0\) , we deduce exact sequences

\[
0 \longrightarrow \operatorname{Ext} _ {\mathrm{A}} ^ {i} (\mathrm{N}, \mathrm{M}) \longrightarrow \operatorname{Ext} _ {\mathrm{A}} ^ {i} (\mathrm{N}, \mathrm{M} _ {1}) \xrightarrow {\varphi^ {i}} \operatorname{Ext} _ {\mathrm{A}} ^ {i + 1} (\mathrm{N}, \mathrm{M}) \longrightarrow 0.
\]

Since \(\operatorname{Ext}_{\mathbf{A}}^{i}(\mathbf{N}, \mathbf{M}_{1}) = 0\) for \(i < n - 1\) , we deduce that \(\operatorname{Ext}_{\mathbf{A}}^{i + 1}(\mathbf{N}, \mathbf{M}) = 0\) for \(i < n - 1\) , that is, \(i + 1 < n\) ; it follows that \(\varphi^{i}\) is bijective for \(i < n\) . Since \(\psi^{i}(\alpha) = \theta_{x} \circ \alpha = \theta_{x_{1}} \circ \theta_{x'} \circ \alpha = \varphi^{i - 1} \circ \overline{\psi}^{i}(\alpha)\) for \(\alpha \in \operatorname{Ext}_{\mathbf{A}}^{i - n}(\mathbf{N}, \mathbf{M}/(\mathbf{x})\mathbf{M})\) , \(\psi^{i}\) is indeed bijective for \(i \leqslant n\) .

\backmatter
